\section{Virtual Ledgers, Strategies, and the Total Return Swap}
\label{ch:virtual}

This chapter discharges constitution clauses C-10.1--C-10.5 (the virtual ledger, the strategy contract, the two disciplines, exact slippage, and the total return swap).

Nothing in this architecture is specific to real holdings. The machine of
Chapters~2 through~7 --- units, wallets, balances, an immutable log, and the
three machines that map events into transactions and fold transactions into
state --- makes no reference to whether the positions it records will ever
settle. A \emph{virtual ledger} is a complete instance of that same machine
whose holdings are notional: it settles nothing, and none of its wallets ever
reaches a settlement instruction or the balance sheet. Everything else is
identical. It has its own units and wallets, its own append-only hash-chained
log, its own Event Monitor, Events Executor, and Transaction Executor, and it is
subject to the same eight commitments as a true ledger --- above all to full
reproducibility, so that any party holding the record recomputes its every
number to the minor unit, and to time travel, so that every past state is
reconstructible from that record alone.

The word \emph{virtual} carries exactly one meaning throughout this
specification, and no other: on the far side of the settlement boundary. A
virtual ledger is not an approximation, a draft, or a lower grade of record. It
is held to the identical standard of correctness; it differs from a true ledger
in one respect only --- its balances are claims about a hypothetical world, so
no move it records is ever projected into a payment. This is the same boundary
the settlement interface draws for a true ledger, read from the other side: the
true ledger says what settles, the virtual ledger says what nothing settles.

\subsection{The strategy as a smart contract}

The motivating case is a quantitative investment strategy: a rulebook that,
reading market observations, decides a hypothetical portfolio and rebalances it
on a schedule. The architecture already has the parts to record this without a
new primitive.

A deterministic strategy is a smart contract, and its rulebook is the declared
terms of a \emph{strategy unit} --- a non-valued unit (Chapter~3), whose price is
undefined, not zero, because pricing the governance shell beside the portfolio it
governs would count one economic reality twice. The hypothetical portfolio the
rulebook manages is an ordinary wallet in a virtual ledger. Each rebalancing is a
recorded transaction: the rebalancing date is a declared watch on the strategy
unit; the Event Monitor emits the event; the strategy contract fires, reads the
stamped market observations and the portfolio as the virtual ledger then stands,
and proposes the moves the rules require; the Transaction Executor admits them
onto the virtual log. The contract never waits and never remembers --- whatever
it must carry from one rebalancing to the next it writes into the three homes,
exactly as every other contract does.

The strategy's \emph{level} is that wallet's Net Asset Value: the same NAV
projection of Chapter~7, the same fold of the same log against a price vector,
computed on demand and never stored. Because the strategy ledger is replayable
and the rebalancing contract is a pure function of recorded rules and recorded
observations, any party holding the recorded rulebook and the recorded
observations recomputes the level exactly --- which is the whole content of the
claim that the level is well defined. There is no privileged copy of it inside
the strategy provider.

An index restatement follows the auditability commitment without exception. When
the level of a published index is corrected --- a late-arriving observation, a
provider's amended fixing --- the virtual ledger does not rewrite its log. It
records a compensating entry that names the transaction it supersedes, so that
both the original computation and its correction survive on the record. A
restatement is an entry, never a rewrite; the historical level as first published
and the historical level as corrected are both reconstructible, and which one a
consumer read on which date is itself a fact of the record.

\subsection{Two disciplines}

Two disciplines make the arrangement safe, and both are checkable.

\emph{The bridge between ledgers is an observation, never a move.} A move
transfers a positive quantity within one ledger; conservation is stated per
ledger, per unit, per coordinate, and it is closed by the virtual wallets that
represent every outside party (Chapter~3). No move therefore spans two ledgers,
and none can: a move from a wallet of \texttt{V-1} to a wallet of the true ledger
has no paired leg that conserves in either, so the Transaction Executor of
neither can admit it. The only thing that ever crosses is a number --- a stamped
level, admitted into the receiving ledger through the observation door of
Chapter~8, with its provenance recorded like any other observation. The two
ledgers touch at exactly one point, and at that point mass does not flow; a fact
is copied.

\emph{The wall between models and the ledger holds level by level.} A ledger
records the outputs of its models and never runs one (Chapter~16); the true
ledger does not execute the strategy, it consumes the level as an observation.
This is not a property of the top ledger alone. Virtual ledgers may nest --- a
strategy over strategies --- and at each level the same wall stands: each ledger
relates to the one it references only through stamped observations admitted at
its boundary, never by reaching inside to run the other's machinery. Because the
wall is re-erected at every level, no amount of nesting lets a model's execution
leak into an accounting record; each ledger remains a pure fold of its own log.

A short type sketch makes both disciplines mechanical:

\begin{lstlisting}
-- a virtual ledger has the very same type as a true one;
-- only the settlement projection distinguishes them.
type VirtualLedger = Ledger

-- the sole thing that crosses the boundary is a stamped level:
observeLevel :: VirtualLedger -> Wallet -> Observation Money

-- a move lives inside one ledger and has no cross-ledger form;
-- none can be given a conserving type, so none can be written.
\end{lstlisting}

\subsection{Real execution: orders, benchmark, and slippage}

When the strategy's trades are meant to be real, the strategy's output toward the
world is an \emph{order} --- an instruction, not a fact. An order records an
intention to trade; it changes no balance, because nothing has happened yet. What
happens is external: the market fills the order, wholly, partly, or not at all,
at a price the strategy did not set. Those fills return as external events,
captured at the boundary, mapped by the responsible contracts, and folded into a
real wallet of the true ledger --- the ordinary path of any executed trade.

The virtual ledger now becomes the benchmark. It holds what the rules say the
portfolio should be; the real wallet holds what execution achieved. Both are NAV
projections, each a fold of its own log against the same prices. Their difference
is the slippage:
\[
  \text{slippage} \;=\; \mathrm{NAV}(\text{benchmark wallet in } V) \;-\;
                        \mathrm{NAV}(\text{real wallet in the true ledger}).
\]
It is the exact difference of two folds, computed from two records, never an
estimate --- every cost that separates the idealised portfolio from the realised
one is captured here in full, because both sides are recorded to the minor unit
and neither is a model output.

\subsection{The total return swap}

A total return swap (TRS) is then one definition and no new machinery: its
reference leg is the NAV of a designated wallet or virtual ledger. The swap
itself is an ordinary unit of the \emph{true} ledger. Its two counterparties hold
wallets there; its smart contract, at each reset, reads the stamped NAV
observation and emits the reset and financing moves. The reference NAV is not
computed by the TRS --- it enters through the door of Chapter~8, exactly as the
strategy level does, and the wall between the two ledgers is the discipline above.

\paragraph{Episode TRS-1.} \textsc{Trs-1} has notional \USD{10{,}000{,}000}. Its
reference leg is the NAV of wallet \texttt{W-QIS} in virtual ledger \texttt{V-1},
managed by the declared rulebook of the non-valued strategy unit \texttt{S-QIS}.
Its financing leg is fixed at 4\% per annum, reset quarterly ($\Delta t = 0.25$).
At inception the reference NAV observation is \USD{10{,}000{,}000}, and the
reference level is set to that figure.

\emph{The trigger.} The first quarterly reset date arrives. \texttt{V-1} has been
running: \texttt{S-QIS}'s rebalancings are recorded transactions on its log, and
its level is the replayable NAV of \texttt{W-QIS}. On the reset date that level is
stamped --- \USD{10{,}300{,}000} --- and admitted into the true ledger as an
observation with provenance (Chapter~8). The Event Monitor emits the reset event
for \textsc{Trs-1}.

\emph{The contract that catches it.} The Events Executor fires \textsc{Trs-1}'s
smart contract, handing it the reset event and the current projection of the true
ledger. The contract reads the stamped reference NAV and the reference level
carried since inception, and computes one net move.

\emph{The transaction it produces.}

\begin{center}
\begin{tabular}{clr}
\toprule
\# & Component & Quantity (USD) \\
\midrule
1 & Total-return leg: stamped NAV $10{,}300{,}000 - $ reference $10{,}000{,}000$ & $300{,}000$ \\
2 & Financing leg: $10{,}000{,}000 \times 4\% \times 0.25$                       & $100{,}000$ \\
3 & Net move --- \texttt{owned(USD)}, payer $\rightarrow$ receiver               & $200{,}000$ \\
4 & Reference level reset (UnitStatus): $10{,}000{,}000 \rightarrow 10{,}300{,}000$ & --- \\
\bottomrule
\end{tabular}
\end{center}

The total-return leg is the appreciation of the reference over the period,
\USD{300{,}000}; the financing leg accrues on the period's opening reference,
\USD{100{,}000}; their difference nets to a single move of \USD{200{,}000} of
\texttt{owned(USD)} from the payer's wallet to the receiver's wallet. In the same
transaction the contract writes the new reference level, \USD{10{,}300{,}000},
into the swap's UnitStatus, arming the next period against it. One reset, one net
move, one state change.

\emph{What the door checks.} The Transaction Executor admits the transaction on
its structure alone. The \USD{200{,}000} move is a paired-leg
\texttt{owned(USD)} transfer, so conservation holds by construction. The
reference-level write has a single legal writer --- the swap's UnitStatus --- and
lands in its one home. The proposal carries a cause-derived identifier keyed to
the reset observation, so a retried or duplicated reset commits once and no
second time. The stamped NAV that the contract read is itself on the record, with
its provenance, so the whole reset is reproducible: re-running the map on the
recorded observation reproduces the identical \USD{200{,}000}.

\emph{The design point.} \emph{The reset moves cash in the true ledger, and the
strategy that produced the number never touched it: the only thing that crossed
the settlement boundary was a stamped observation, and a move never can.}

\subsection{The simulated ledger}\label{sec:simulated}

Concrete case first. At the close of a recorded trading day the true ledger stands in a
definite state: its log has a last position. Take that position as a \textbf{branch point},
copy the log up to it, keep every contract and every machine, and continue it --- not with
the day's real events, but with a generated stream of clock ticks and boundary observations.
The result is a \textbf{simulated ledger}: a virtual ledger (this chapter) whose log extends
a recorded prefix of another ledger's log, driven onward by generated events. Branch
\texttt{V-SIM} at that close and drive OT-1 --- the one-touch on OMEGA struck at barrier
$80.00$ (Chapter~\ref{ch:objects}) --- down $1{,}000$ generated OMEGA paths. Each path is a
sequence of generated official-close observations admitted through the observation door
(Chapter~\ref{ch:marketdata}); on each, OT-1's contract fires unchanged, and a path whose
close touches $80.00$ pays the declared \USD{1{,}000{,}000} while one that never does pays
nothing. The price estimate is the mean payout across the thousand paths. No new machinery
produced it: the same door, the same contract, a generated stream in place of the real one.

\emph{The generator is outside the ledger.} The stream of ticks and observations is produced
by a \textbf{generator} --- a model, and a model lives outside the fold (the wall between
models and the ledger, Chapter~16). The generator's one
non-record input is the \textbf{seed}. The seed is recorded, so a \textbf{path} is a
deterministic function of the branch point, the generator version, and the seed: fix those
three and the path replays bit for bit. Nothing else enters. This is the confinement the
constitution's simulability commitment demands (C-2.8) --- one non-record input, recorded,
so every path replays exactly.

\emph{No simulation-only variant exists.} This is the section's discipline, and it is
absolute. There is no simulation build of any contract, and none of the Event Monitor, the
Events Executor, or the Transaction Executor. The generated stream is admitted through the
one door that admits real events and folded by the same contracts into the same three homes
(Chapter~\ref{ch:homes}). A contract that behaved one way in production and another under
simulation would be an untested second implementation --- the exact divergence the
one-record architecture abolishes. The simulated ledger is held to the same eight
commitments as any ledger; it differs from the true ledger in one respect only, that its
driving events were generated rather than observed.

\emph{Results cross back as observations.} A simulated ledger settles nothing, exactly as
any virtual ledger. A number computed on it --- the price of OT-1, a stress loss, a
sensitivity --- re-enters the true ledger only as an observation with provenance, through
the observation door of Chapter~\ref{ch:marketdata}, by the same moveless
observation-recording transaction that admits any stamped level. The provenance a simulated
result carries names its branch point, the generator version, the seeds, and the path count,
so the figure is reproducible from the record --- and no move ever crosses from the
simulated ledger to the true one.

\emph{One engine, two measures.} Chapter~\ref{ch:testability}'s generator regime is this same
engine. Its generated products, events, and histories are simulated paths branched from
recorded states and folded by the same contracts and the same door; only their purpose
differs. A path driven to price OT-1 runs to ask what the instrument is worth; a path driven
by the adversarial scheduler runs to ask whether an invariant can be broken. The machinery
is identical --- one generator universe, one door, one fold --- and the two uses are two
measures over the same paths. Refutation and pricing are one computation aimed at different
questions.

\begin{lstlisting}
-- a branch point is a recorded prefix of a ledger's log:
branch        :: Ledger -> LogPrefix -> BranchPoint
-- a path is a deterministic function of (branch point, generator version, seed),
-- run through the same contracts and the same door as production:
simulate      :: BranchPoint -> GeneratorVersion -> Seed -> Path
-- a number computed on a path re-enters the true ledger only as a stamped
-- observation with provenance -- never as a move (Chapter 8):
observeResult :: Path -> Observation Money
\end{lstlisting}
